% ============================================================
% Section 184: δ势束缚态在电场中的电离概率
% ============================================================
\section{量子力学：$\delta$势束缚态在电场中的电离概率}
\label{sec:delta-ionization-probability}

\begin{modelbox}[题目：$\delta$势束缚态在电场中的电离概率]
电荷为 $q$ 的粒子受位于原点的 $\delta$ 势 $V(x)=-A\delta(x)$ 束缚。$t=0$ 到 $t=\tau$ 内受到沿 $x$ 方向的匀强电场 $\mathcal{E}_0$ 作用。求 $t>\tau$ 时粒子处于能量范围 $E_k\to E_k+dE_k$ 的非束缚态的概率 $P(E_k)dE_k$。
\end{modelbox}

\begin{tipbox}[考点快速分析]
\begin{itemize}
    \item \textbf{束缚态}：$\psi_0(x)=\sqrt{\kappa}\,e^{-\kappa|x|}$，$\kappa=mA/\hbar^2$，$E_0=-mA^2/(2\hbar^2)$
    \item \textbf{自由态箱归一化}：$\psi_k(x)=e^{ikx}/\sqrt{L}$，态密度 $D(E_k)=\dfrac{L}{\pi\hbar}\sqrt{\dfrac{m}{2E_k}}$
    \item \textbf{微扰矩阵元}：$\langle k|H_1|0\rangle\propto\dfrac{k}{(k^2+\kappa^2)^2}$（奇函数）
    \item \textbf{时间积分}：$\sin^2(\omega\tau/2)/\omega^2$ 衍射因子
\end{itemize}
\end{tipbox}

\begin{derivationbox}[标准解题过程]
\textbf{1. 束缚态}

$\psi_0(x)=\sqrt{\kappa}\,e^{-\kappa|x|}$，$\kappa=mA/\hbar^2$，$E_0=-\hbar^2\kappa^2/(2m)=-mA^2/(2\hbar^2)$。

\textbf{2. 自由粒子态与态密度}

箱归一化 $L$（周期性边界）：$\psi_k(x)=\dfrac{1}{\sqrt{L}}e^{ikx}$，$k=2\pi n/L$。

态密度（含 $\pm k$ 二重简并）：
\[
D(E_k)=\frac{L}{\pi\hbar}\sqrt{\frac{m}{2E_k}}.
\]

\textbf{3. 微扰矩阵元}

$H_1(t)=-q\mathcal{E}_0x$（$0\le t\le\tau$）。
\[
\langle k|H_1|0\rangle=-q\mathcal{E}_0\sqrt{\frac{\kappa}{L}}\int_{-\infty}^{\infty}xe^{-ikx}e^{-\kappa|x|}dx=\frac{4iq\mathcal{E}_0\kappa^{3/2}}{\sqrt{L}}\frac{k}{(k^2+\kappa^2)^2}.
\]

矩阵元对 $\pm k$ 为奇函数，模方相同。

\textbf{4. 跃迁概率}

含时微扰论：
\[
P_{0\to k}=\frac{1}{\hbar^2}|\langle k|H_1|0\rangle|^2\frac{4\sin^2(\omega_{k0}\tau/2)}{\omega_{k0}^2},
\]
其中 $\omega_{k0}=(E_k-E_0)/\hbar=\dfrac{\hbar}{2m}(k^2+\kappa^2)$。

对能量区间 $dE_k$ 内的所有末态求和（含 $\pm k$ 二重简并）：
\begin{equation}
\boxed{P(E_k)dE_k=\frac{256m^3q^2\mathcal{E}_0^2\kappa^3}{\pi\hbar^6}\frac{\sqrt{2mE_k}}{\bigl(\kappa^2+2mE_k/\hbar^2\bigr)^6}\sin^2\!\Bigl(\frac{\hbar\tau}{4m}\bigl(\kappa^2+\frac{2mE_k}{\hbar^2}\bigr)\Bigr)dE_k.}
\end{equation}

物理特征：
\begin{itemize}
    \item $\sin^2$ 衍射因子使能量分布呈现振荡，中心峰值在 $\omega_{k0}\approx0$（低能）附近；
    \item 随 $E_k$ 增大，$P(E_k)\propto E_k^{-5/2}$ 快速衰减；
    \item 电场越强（$\mathcal{E}_0$ 大）或作用时间越长（$\tau$ 大），总电离概率越大。
\end{itemize}
\end{derivationbox}

\begin{formulabox}[答案汇总]
\begin{center}
\begin{tabular}{ll}
\toprule
\textbf{结论} & \textbf{结果} \\
\midrule
束缚态波函数 & $\psi_0=\sqrt{\kappa}\,e^{-\kappa|x|}$，$\kappa=mA/\hbar^2$ \\
自由态态密度 & $D(E_k)=\dfrac{L}{\pi\hbar}\sqrt{\dfrac{m}{2E_k}}$ \\
跃迁矩阵元 & $\langle k|H_1|0\rangle\propto\dfrac{k}{(k^2+\kappa^2)^2}$ \\
能量分布 & $P(E_k)\propto\dfrac{\sqrt{E_k}}{(\kappa^2+2mE_k/\hbar^2)^6}\sin^2(\cdots)$ \\
\bottomrule
\end{tabular}
\end{center}
\end{formulabox}

\begin{insightbox}[物理图像]
\begin{itemize}
    \item 本题是含时微扰论处理电离过程的典型范例。$\delta$ 势的简单性使解析计算成为可能，而物理图像具有普遍性。
    \item $\sin^2(\omega\tau/2)/\omega^2$ 因子是有限时间微扰的普遍特征：当 $\tau\to\infty$ 时趋于 $\pi\tau\delta(\omega)$（Fermi黄金定则）；有限 $\tau$ 下能量守恒放宽为 $\Delta E\sim\hbar/\tau$。
    \item 该模型可推广到理解光电效应、隧穿电离和强场物理中的 Above-Threshold Ionization (ATI)。
\end{itemize}
\end{insightbox}
